\section{Composition controls and a diagnosis of the 60:40 comparison}
\label{sec:composition-diagnosis-next}
The composition experiment changes the fractions while retaining the
aggregate input matrices. Write $F=\operatorname{diag}(f_E,f_I)$,
$W=GF$, and $B=J^0F$. The specified path is
\begin{equation}
 G(f)=G_{80}\operatorname{diag}(.8/f_E,.2/f_I),\qquad
 J^0(f)=J^0_{80}\operatorname{diag}(.8/f_E,.2/f_I).
\end{equation}
Consequently $W$ and $B$ are constant. The leading conditional one-site law
is unchanged on this path. The Bernoulli probabilities
$p_{ab}=(J^0_{ab})^2/[(J^0_{ab})^2+G_{ab}]$ change with composition,
so higher connectivity cumulants and finite-network fluctuations need not
be unchanged. This distinction prevents an inference that finite networks
must share identical local statistics.

\subsection{Population geometry with an unambiguous normalization}
For the stationary temporally centered rate covariance $C$, define
\begin{align}
 \nu_a&=\frac1{N_a}\sum_{i\in a}C_{ii},\\
 \mathcal E_{ab}&=\frac1N\sum_{i\in a,j\in b}C_{ij}^2,\\
 \Xi_{ab}&=\frac{N}{N_aN_b}\sum_{i\in a,j\in b,
                    (a,i)\ne(b,j)}C_{ij}^2.
\end{align}
The unequal-population exclusion is vacuous. The two ordered cross blocks
are both included: $\mathcal E_{EI}=\mathcal E_{IE}$ at zero lag and their
sum is twice the rectangular contribution. With
$V_{2,a}=N_a^{-1}\sum_{i\in a}C_{ii}^2$,
\begin{equation}
 \nu=\sum_af_a\nu_a,\qquad
 \mathcal E=\sum_{ab}\mathcal E_{ab}
 =\sum_af_a V_{2,a}+\sum_{ab}f_af_b\Xi_{ab},\qquad
 \frac DN=\frac{\nu^2}{\mathcal E}.
\end{equation}
The within-population quantity is instead
$D_a/N_a=\nu_a^2/(V_{2,a}+f_a\Xi_{aa})$.
An E--I rectangular cross-covariance block does not have this PR definition.
The conditional leading theory predicts these population energies, not only
their contraction into a single dimension.

\subsection{An exact accounting of ingredient errors}
For any positive measured $\nu_N,\mathcal E_N$ and theoretical
$\nu_0,\mathcal E_0$, no small-error expansion is needed for
\begin{equation}
 \log\frac{d_N}{d_0}
 =2\log\frac{\nu_N}{\nu_0}
  -\log\frac{\mathcal E_N}{\mathcal E_0}.
 \label{eq:exact-log-error-next}
\end{equation}
For ensemble means the same equation applies to
$d_{\rm plug}=(\mathbb E\nu_N)^2/\mathbb E\mathcal E_N$.
It does not apply to the mean network ratio without an additional term.
Putting $a=(\mathbb E\nu_N)/(\mathbb E\mathcal E_N)$ and expanding a
nonnegative square gives
\begin{align}
 \mathbb E\frac{(\nu_N-a\mathcal E_N)^2}{\mathcal E_N}
 &=\mathbb E\frac{\nu_N^2}{\mathcal E_N}
       -2a\mathbb E\nu_N+a^2\mathbb E\mathcal E_N\\
 &=\mathbb E d_N-\frac{(\mathbb E\nu_N)^2}
                            {\mathbb E\mathcal E_N}\ge0.
\end{align}
Thus the mean ratio and the ratio of ensemble moments are distinct observables.
They must both be retained rather than selecting whichever agrees better.

In the historical sixteen-network 60:40 cohort with the four-window estimator
at $T=16000$ and $B=500$, $\bar\nu/\nu_0=0.94218$ and
$\bar{\mathcal E}/\mathcal E_0=1.09318$. Equation~\eqref{eq:exact-log-error-next}
then has contributions $-0.11912$ and $-0.08909$ respectively.
The ratio of mean moments is $0.0242746$, whereas the mean network ratio is
$0.0256785$. These are finite-record observations in a historically assembled
cohort, not an infinite-time ensemble theorem. They establish that the
remaining difference cannot be assigned solely to the old own-block
centering artifact or solely to the two original unusually low-PR networks.

\subsection{What the formal first correction can and cannot explain}
Write $\epsilon=N^{-1/2}$ and
\begin{equation}
 \nu=\nu_0+\epsilon\nu_1+O(\epsilon^2),\quad
 \mathcal E=\mathcal E_0+\epsilon\mathcal E_1+O(\epsilon^2).
\end{equation}
Expanding the ratio explicitly,
\begin{align}
 \frac{\nu^2}{\mathcal E}
 &=\frac{\nu_0^2+2\epsilon\nu_0\nu_1+O(\epsilon^2)}
         {\mathcal E_0}
       \left(1-\epsilon\frac{\mathcal E_1}{\mathcal E_0}
             +O(\epsilon^2)\right)\\
 &=d_0\left[1+\epsilon\left(
       \frac{2\nu_1}{\nu_0}-\frac{\mathcal E_1}{\mathcal E_0}
                               \right)\right]+O(\epsilon^2).
\end{align}
For the corrected formal 60:40 coefficients, the $N=1000$ retained shifts
are $+8.806\%$ in $\nu$, $+14.341\%$ in $\mathcal E$, and $+3.271\%$
in $d$. The observed historical $\nu$ shift is negative and the observed
mean ratio shift is negative. The retained correction therefore does not
repair this comparison. Replacing the strict expansion by a ratio of
truncated moments silently retains some second-order terms but omits others;
that change is not evidence of a controlled repair.

\subsection{A diagnostic substitution is not a new prediction}
On a common frequency grid put
\begin{equation}
 T=\operatorname{diag}(A)H(\omega_1)H(\omega_2)W,\quad
 R=(I-T)^{-1},\quad \Sigma=\operatorname{diag}(Q/f),\quad
 \mathcal E=\iint f^\top R\Sigma R^\dagger f.
\end{equation}
The source $Q_a$ is a same-neuron product of spectra. It is not the product
of population-averaged spectra. Three substitutions distinguish errors:
replace only $Q$ and the numerator; replace only $A$; replace both. All use
measured local statistics and are calibrated diagnostics. A pole after the
substitution is a failed diagnostic, not an invitation to clip the gains.
For a small perturbation,
\begin{align}
 (I-T-\delta T)(R+\delta R)&=I+O(\delta T^2),\\
 (I-T)\delta R&=\delta T R+O(\delta T^2),\\
 \delta R&=R\delta T R+O(\delta T^2),\\
 \delta\mathcal E
 &=\iint f^\top R\delta\Sigma R^\dagger f
  +2\operatorname{Re}\iint f^\top R\delta T R\Sigma R^\dagger f
  +O(\delta^2).
\end{align}
The two changes can oppose one another and their correlation is physical.
A gain-only improvement at fixed source is insufficient.
Prior held-out substitutions already show that measured one-site statistics
are not a universal pair-energy repair: the seed-22 $N=1000$ case is close,
whereas seed 11 and both tested $N=2000$ cases retain substantial errors.
Those earlier estimates used a different recording estimator, so they
motivate the new common-protocol comparison rather than supplying its result.

\subsection{An exact test for a population-coherent energy sector}
Let $U_{ia}=1_{i\in a}/\sqrt{N_a}$, $P=UU^\top$, and $Q_\perp=I-P$.
For a covariance matrix $C$ define $A=U^\top CU$ and $Z=CU$.
Using orthogonal projectors, write
\begin{equation}
 C=PCP+PCQ_\perp+Q_\perp CP+Q_\perp CQ_\perp.
\end{equation}
The four blocks are Frobenius orthogonal, whence
\begin{align}
 \|C\|_F^2
 &=\|A\|_F^2+2\|PCQ_\perp\|_F^2
                     +\|Q_\perp CQ_\perp\|_F^2,\\
 \|Z\|_F^2&=\|A\|_F^2+\|Q_\perp CPU\|_F^2,\\
 \mathcal E_{\rm touching\ population}
 &=\frac{2\|Z\|_F^2-\|A\|_F^2}{N},\\
 \mathcal E_\perp
 &=\mathcal E-\mathcal E_{\rm touching\ population}.
\end{align}
The second identity follows from $CU=PCU+Q_\perp CU$ and $U^\top U=I$.
The displayed redundant $P U=U$ makes its dimensions explicit.
For two covariance estimates $C,D$, polarization replaces squared norms by
inner products. Thus the same decomposition applies to cross-window energy
and to the mean-difference-corrected covariance matrices, without estimating
a full eigenspectrum. Finite-record cross sectors may be negative, even
though every stationary self-energy sector is nonnegative.
A large population sector identifies coherent geometry but does not prove a
specific feedback mechanism; a small sector rules out population-uniform
modes as the whole explanation. Nonuniform collective modes can remain in
$Q_\perp$.

\subsection{Static--dynamic variance partition and a conditional susceptibility}
The law of total variance gives, within a population,
\begin{equation}
 \operatorname{Var}_{i,t}x_i(t)=\Delta+q,\qquad
 \operatorname{Var}_{i,t}r_i(t)=\operatorname{Var}_i m_i+\nu.
\end{equation}
Consequently a good check of the marginal input variance alone does not
establish the correct temporal variance. A finite network can put more of
its activity variation into fixed differences between neurons and less into
temporal fluctuations.
There is an exact conditional Gaussian derivative behind this interpretation.
Temporarily hold $h$ and total input variance $V=\Delta+q$ fixed and put
\begin{equation}
 M_2(h,V,\Delta)=\mathbb E[\phi(X_1)\phi(X_2)],\qquad
 \operatorname{Var}X_j=V,\quad\operatorname{Cov}(X_1,X_2)=\Delta.
\end{equation}
Price's theorem and conditional independence given the static label imply
\begin{align}
 \partial_\Delta M_2\big|_{h,V}
 &=\mathbb E[\phi'(X_1)\phi'(X_2)]\\
 &=\mathbb E_z\left[\mathbb E_u\phi'(h+\sqrt\Delta z+u)\right]^2=A,\\
 \nu(h,V,\Delta)&=M_2(h,V,V)-M_2(h,V,\Delta),\\
 \partial_\Delta\nu\big|_{h,V}&=-A<0.
\end{align}
For $\phi=2\Phi$, $s=1+V$ gives explicitly
\begin{equation}
 A=\frac{2}{\pi\sqrt{s^2-\Delta^2}}
       \exp\left[-\frac{h^2}{s+\Delta}\right],\qquad
 \partial_\Delta\log A\big|_{h,V}
 =\frac{\Delta}{s^2-\Delta^2}+\frac{h^2}{(s+\Delta)^2}\ge0.
\end{equation}
Thus a pure repartition at fixed $h,V$ lowers temporal rate variance and
locally increases this gain moment when $\Delta>0$ or $h\ne0$. If the measured gain moves in the other direction,
changes in $h,V$, the static label distribution, or conditional temporal
statistics are also needed. The partition argument alone is not a complete
finite-network explanation.

For a prescribed small residual $\xi$ in the Gaussian static equation
$\Delta=W M_2(\Delta)+\xi$, with $h,V$ held fixed, its linear response is
\begin{equation}
 \delta\Delta=(I-W\operatorname{diag}A)^{-1}\delta\xi.
\end{equation}
At the reference law the largest eigenvalue of
$W\operatorname{diag}A$ is $0.8558355$, the modal amplification is $6.93652$,
and the inverse matrix has two-norm $10.17336$. The associated positive
right eigenvector is $(0.546639,0.837368)^\top$ with unit Euclidean norm.
These values identify a sensitive partition direction in the conditional
functional. They are not the full quenched fluctuation law: the actual
network also changes $V,h$, its temporal kernel and its source residual.

A separate parameter-only check varies the degree dispersion in the full
population boundary-value problem. At $\chi=1$,
\begin{equation}
 \frac{d\Delta}{d\chi}=(2.800885,4.294997)^\top,\qquad
 \frac{dq}{d\chi}=(-0.907054,-1.388953)^\top.
\end{equation}
Central differences with steps $0.01$ and $0.005$ agree to about
$10^{-5}$ relatively. Total variance also increases on this path;
$\Delta+q$ is not fixed. This is a controlled test of the disorder mechanism,
not a fit of an effective degree dispersion to a measured PR.

\subsection{Exact static-source accounting on a fixed network}
Write $J_{ij}=J^0_{ab}/\sqrt N+\delta J_{ij}$ and let $m_j$ be the
measured time-averaged firing rate. Define
$\eta_i=\sum_j\delta J_{ij}m_j$. Within population $a$, the coherent
mean input is common to every neuron, so the stationary static input
variance is $\operatorname{Var}_{i\in a}\eta_i$. Direct expansion gives
\begin{align}
 \Delta_a
 &=\frac1{N_a}\sum_{i\in a}\sum_j\delta J_{ij}^2m_j^2
 +\frac1{N_a}\sum_{i\in a}\sum_{j\ne k}
           \delta J_{ij}\delta J_{ik}m_jm_k
 -\left(\frac1{N_a}\sum_{i\in a}\eta_i\right)^2.
 \label{eq:fixed-static-source-accounting}
\end{align}
This is an identity for fixed connectivity and fixed measured $m$; it does
not use independence of the rates and the connectivity. The leading
Bernoulli closure replaces the first term by
$\sum_bW_{ab}\mathbb E_{j\in b}m_j^2$ and discards the other centered
terms in the appropriate ensemble limit. Comparing these pieces diagnoses
three different changes: the empirical presynaptic rate distribution,
the realized squared couplings weighted by those rates, and the remaining
cross-source/centering contribution. The identity alone predicts none of
their finite-network distributions.
A check of the measured mean state against $Jm+\mathrm{drive}$ accompanies
this calculation, since finite-record endpoint terms and sampling a
midpoint integrator can leave a small difference.

For Bernoulli nonzero signed weight $t=j/p$, write $j=J^0_{ab}$ and
$G=G_{ab}$. Elementary substitution gives
\begin{align}
 p&=\frac{j^2}{j^2+G},\quad t=j+G/j,\\
 T_3&=t^3p(1-p)(1-2p)=G(G-j^2)/j,\\
 T_4&=t^4p(1-p)[1-6p(1-p)]
        =G\left(j^2-4G+\frac{G^2}{j^2}\right).
\end{align}
On the fixed-aggregate path $j=B_{ab}/f_b$, $G=W_{ab}/f_b$,
\begin{equation}
 f_bT_{3,ab}=\frac{W_{ab}^2}{B_{ab}}
                   -\frac{W_{ab}B_{ab}}{f_b},\qquad
 f_bT_{4,ab}=\frac{W_{ab}B_{ab}^2}{f_b^2}
             -\frac{4W_{ab}^2}{f_b}+\frac{W_{ab}^3}{B_{ab}^2}.
\end{equation}
Thus holding $W,B$ fixed does not hold the higher finite-size source
cumulants fixed. This states the distinction without asserting a completed
second-order correction.
The same decomposition also keeps the finite coherent balance feedback visible.
For external input $-\sqrt N Bm^*+h^{\rm ext}$, average the stationary
mean equation first over time and then over a population:
\begin{equation}
 \bar\mu=h^{\rm ext}+\sqrt N B(\bar m-m^*)+\bar\eta,
 \qquad
 \bar m-m^*=N^{-1/2}B^{-1}(\bar\mu-h^{\rm ext}-\bar\eta).
\end{equation}
The centered recurrent input $\bar\eta$ cannot in general be dropped in a
finite network. In a finite record the equation also has an endpoint
contribution from the mean derivative. The selected-network source check
retains this term as a measured residual rather than treating the inferred
balance identity as a new parameter-only prediction.

\subsection{Population dimensions and a normalized cross-population energy}
The population PR is a square-block quantity:
\begin{equation}
 d_a\equiv\frac{D_a}{N_a}
   =\frac{f_a\nu_a^2}{\mathcal E_{aa}},\qquad
 D_a=Nf_a d_a.
\end{equation}
Reporting both $d_a$ and $D_a$ distinguishes the fraction of population size
from the absolute effective number of dimensions. A smaller population can
have a larger dimension fraction while exploring fewer absolute dimensions.
The full-network PR need not equal the sum of the population PRs.

A separate normalized observable uses the cross-population covariance energy:
\begin{equation}
 \rho_{\rm geom}=
 \frac{\|C_{EI}\|_F^2}{\|C_{EE}\|_F\|C_{II}\|_F}
 =\frac{\mathcal E_{EI}}
        {\sqrt{\mathcal E_{EE}\mathcal E_{II}}}.
\end{equation}
To see its meaning, define temporal Gram kernels
$K_a(t,s)=\sum_{i\in a}\delta r_i(t)\delta r_i(s)$.
Averages over two independent stationary times give
\begin{equation}
 \langle K_aK_b\rangle_{t,s}=\sum_{i\in a,j\in b}C_{ij}^2,
 \qquad
 \rho_{\rm geom}=\frac{\langle K_EK_I\rangle}
                        {\|K_E\|\|K_I\|}\in[0,1],
\end{equation}
where the last inequality is Cauchy--Schwarz. The normalization removes a
common amplitude from each population. It tests their temporal covariance
geometry, not a PR of a rectangular matrix or a task-based communication
capacity. Cross-window finite-record estimates need not themselves respect
all stationary positive-semidefinite inequalities.
